\section{Đặc tính kiến trúc trong IoT}

\subsection{Khái niệm đặc tính kiến trúc}

Trong kỹ nghệ phần mềm, quá trình thiết kế hệ thống đòi hỏi sự phân định rõ ràng giữa hai khía cạnh yêu cầu cốt lõi \cite{richards2020}:

\begin{itemize}
    \item \textbf{Functional Requirements:} Xác định các hành vi, quy trình nghiệp vụ và luồng dữ liệu mà hệ thống phải thực thi. Khía cạnh này giải quyết bài toán ứng dụng (ví dụ: thu thập và lưu trữ dữ liệu nhiệt độ theo chu kỳ).
    \item \textbf{Architectural Characteristics:} Còn được gọi là thuộc tính chất lượng hoặc yêu cầu phi chức năng, xác định các tiêu chuẩn vận hành mà hệ thống phải đạt được \cite{richards2020}. Khía cạnh này định hình cấu trúc nền tảng (ví dụ: khả năng duy trì độ trễ dưới 100ms khi xử lý 1 triệu kết nối đồng thời).
\end{itemize}

Việc phân lập này mang tính quyết định, bởi lẽ logic nghiệp vụ có thể được tái cấu trúc trong quá trình phát triển, nhưng kiến trúc nền tảng thì rất khó thay đổi sau khi hệ thống đã đi vào vận hành. Một quy trình thiết kế chuẩn mực phải tuân thủ trình tự: \textbf{Yêu cầu nghiệp vụ $\rightarrow$ Đặc tính kiến trúc $\rightarrow$ Mẫu kiến trúc $\rightarrow$ Công nghệ triển khai} \cite{richards2020}.

\subsection{Các đặc tính cốt lõi của hệ thống IoT}

Khác với phần mềm truyền thống, hệ thống IoT phải tương tác trực tiếp với phần cứng, do đó thiết kế kiến trúc phải đối phó với những giới hạn khắt khe hơn từ thế giới thực \cite{lea2018}. Dưới đây là 8 đặc tính cốt lõi mà bất cứ hệ thống IoT nào cũng cần phải có:

\subsubsection*{1. Scalability}
\begin{itemize}
    \item \textbf{Định nghĩa:} Năng lực của hệ thống trong việc duy trì hoặc gia tăng mức độ hiệu năng khi khối lượng công việc được mở rộng \cite{richards2020}.
    \item \textbf{Góc độ IoT đặc thù:} Quá trình mở rộng trong IoT diễn ra trên hai trục độc lập: số lượng kết nối duy trì trạng thái và thông lượng dữ liệu thu nhận.
    \item \textbf{Ví dụ thực tiễn:} Trong một hệ thống IIoT (Industrial IoT), kiến trúc có thể sử dụng Apache Kafka làm lớp đệm phân tán trước khi định tuyến dữ liệu đến các cụm xử lý luồng như Apache Spark \cite{databricks_johndeere}.
\end{itemize}

\subsubsection*{2. Reliability}
\begin{itemize}
    \item \textbf{Định nghĩa:} Khả năng hệ thống duy trì các chức năng được chỉ định trong những điều kiện nhất định và trong một khoảng thời gian xác định \cite{richards2020}.
    \item \textbf{Góc độ IoT đặc thù:} Đặc tính này bao hàm khả năng "tự chủ ngoại tuyến" tại các điểm biên. Thiết bị phải có khả năng tiếp tục vận hành các quy trình thiết yếu ngay cả khi mất kết nối mạng WAN với đám mây \cite{lea2018}.
    \item \textbf{Ví dụ thực tiễn:} Một Edge Gateway tại trạm biến áp tiếp tục đánh giá dữ liệu cảm biến dựa trên công cụ quy tắc cục bộ để kích hoạt rơ-le ngắt điện khẩn cấp khi phát hiện quá tải \cite{lea2018}.
\end{itemize}

\subsubsection*{3. Performance}
\begin{itemize}
    \item \textbf{Định nghĩa:} Mức độ đáp ứng của hệ thống về mặt thời gian và thông lượng khi thực thi các tác vụ \cite{richards2020}.
    \item \textbf{Góc độ IoT đặc thù:} Do tương tác trực tiếp với các quá trình vật lý, yêu cầu về hiệu năng thường đi kèm với các ràng buộc thời gian thực và thời hạn cứng.
    \item \textbf{Ví dụ thực tiễn:} Trong điều khiển robot phẫu thuật từ xa, độ trễ tín hiệu lệnh phải được đảm bảo nghiêm ngặt dưới ngưỡng 50ms để ngăn ngừa rủi ro vận hành vật lý \cite{lea2018}.
\end{itemize}

\subsubsection*{4. Security}
\begin{itemize}
    \item \textbf{Định nghĩa:} Năng lực bảo vệ dữ liệu và hệ thống khỏi sự truy cập, sửa đổi, hoặc phá hoại trái phép.
    \item \textbf{Góc độ IoT đặc thù:} Bề mặt tấn công mở rộng ra môi trường vật lý. Các thiết bị đầu cuối thường đặt tại các vị trí không được kiểm soát an ninh \cite{lea2018}.
    \item \textbf{Ví dụ thực tiễn:} Tích hợp các module phần cứng bảo mật (như TPM) để lưu trữ chứng chỉ khóa công khai (mTLS), đảm bảo không thể trích xuất khóa riêng tư để giả mạo danh tính trên mạng lưới \cite{lea2018}.
\end{itemize}

\subsubsection*{5. Maintainability}
\begin{itemize}
    \item \textbf{Định nghĩa:} Mức độ hiệu quả và dễ dàng khi tiến hành sửa đổi, nâng cấp phần mềm hoặc khắc phục sự cố \cite{richards2020}.
    \item \textbf{Góc độ IoT đặc thù:} Chu kỳ sống của phần cứng kéo dài từ 10 đến 15 năm tại các môi trường khó tiếp cận, đòi hỏi cơ chế cập nhật OTA (Over-The-Air) phải có độ tin cậy tuyệt đối \cite{lea2018}.
    \item \textbf{Ví dụ thực tiễn:} Triển khai kiến trúc bộ nhớ phân vùng kép trên vi điều khiển, cho phép thiết bị tự động khởi động lại bằng phân vùng chứa firmware cũ nếu bản cập nhật phát sinh lỗi \cite{lea2018}.
\end{itemize}

\subsubsection*{6. Interoperability}
\begin{itemize}
    \item \textbf{Định nghĩa:} Năng lực trao đổi và sử dụng thông tin giữa các hệ thống hoặc thành phần phần mềm có kiến trúc khác biệt.
    \item \textbf{Góc độ IoT đặc thù:} Sự phân mảnh nghiêm trọng về tiêu chuẩn công nghiệp đòi hỏi hệ thống phải đồng thời xử lý đa dạng các giao thức truyền thông như: MQTT, CoAP, Modbus, LoRaWAN \cite{lea2018}.
    \item \textbf{Ví dụ thực tiễn:} Sử dụng một Edge Gateway đóng vai trò bộ chuyển đổi giao thức, dịch các bản tin từ Modbus RTU và Profinet sang định dạng chuẩn hóa trước khi truyền qua MQTT tới máy chủ \cite{aws_siemens2019}.
\end{itemize}

\subsubsection*{7. Observability}
\begin{itemize}
    \item \textbf{Định nghĩa:} Mức độ mà trạng thái nội bộ của một hệ thống có thể được suy luận thông qua các dữ liệu đầu ra bên ngoài của nó \cite{garrison2017}.
    \item \textbf{Góc độ IoT đặc thù:} Khả năng truy cập trực tiếp (ví dụ: SSH) vào các thiết bị vi điều khiển bị hạn chế. Hệ thống telemetry phải được thiết kế sẵn để xuất log và metrics hữu ích.
    \item \textbf{Ví dụ thực tiễn:} Hệ thống cảnh báo không chỉ dựa trên trạng thái mất kết nối mà còn tương quan với dữ liệu viễn trắc trước đó (ví dụ: đường cong suy giảm điện áp pin) để phân loại nguyên nhân ngoại tuyến \cite{ms_rollsroyce}.
\end{itemize}

\subsubsection*{8. Energy Efficiency}
\begin{itemize}
    \item \textbf{Định nghĩa:} Tỷ lệ giữa hiệu suất hoạt động của hệ thống so với mức tiêu hao năng lượng.
    \item \textbf{Góc độ IoT đặc thù:} Các thiết bị đầu cuối vận hành bằng pin đòi hỏi kiến trúc truyền thông phải tối ưu hóa chu kỳ hoạt động (duty cycle) của bộ vi xử lý và module vô tuyến \cite{lea2018}.
    \item \textbf{Ví dụ thực tiễn:} Thay thế mô hình giao tiếp đồng bộ dựa trên HTTP bằng giao thức LPWAN (như LoRaWAN), cho phép thiết bị duy trì trạng thái ngủ sâu và chỉ thức giấc theo chu kỳ \cite{lea2018}.
\end{itemize}

\subsection{Sự đánh đổi giữa các đặc tính kiến trúc}

Việc tối ưu hóa đồng thời toàn bộ các đặc tính kiến trúc là bất khả thi về mặt toán học và vật lý. Kiến trúc sư phần mềm bắt buộc phải thực hiện các quyết định đánh đổi chiến lược dựa trên độ ưu tiên của nghiệp vụ \cite{richards2020}:

\begin{itemize}
    \item \textbf{Security vs. Performance:} Quá trình thiết lập phiên bảo mật (ví dụ: TLS handshake) và chi phí tính toán cho các thuật toán mã hóa có thể làm tăng độ trễ mạng. \textit{Phương pháp giảm thiểu:} Ứng dụng tăng tốc phần cứng \cite{lea2018}.
    
    \item \textbf{Reliability vs. Performance:} Trong giao thức MQTT, mức chất lượng dịch vụ cao nhất (QoS 2) đòi hỏi quy trình xác nhận 4 bước, dẫn đến giảm sút nghiêm trọng thông lượng mạng \cite{lea2018}. \textit{Phương pháp giảm thiểu:} Áp dụng định tuyến phân lớp (QoS 0/1 cho telemetry, QoS 2 cho lệnh điều khiển).
    
    \item \textbf{Scalability vs. Maintainability:} Việc phân rã hệ thống thành Microservices giải quyết tốt bài toán phân tải, nhưng làm tăng theo cấp số nhân độ phức tạp vận hành. \textit{Phương pháp giảm thiểu:} Tích hợp công cụ điều phối container (Kubernetes) và quy trình GitOps \cite{garrison2017}.
\end{itemize}

\subsection{Đặc tính kiến trúc dẫn dắt việc lựa chọn mẫu kiến trúc}

Quá trình định lượng và xếp hạng các đặc tính ưu tiên sẽ hình thành cơ sở logic để xác định mẫu kiến trúc chủ đạo \cite{richards2020}:

\begin{itemize}
    \item \textbf{Ưu tiên Performance (Độ trễ < 100ms) $\rightarrow$ Mẫu Edge-Cloud Hybrid:} Kiến trúc phân tán năng lực tính toán về phía Edge Computing là bắt buộc để triệt tiêu độ trễ đường truyền \cite{lea2018}.
    \item \textbf{Ưu tiên Scalability (> 100,000 thiết bị) $\rightarrow$ Microservices kết hợp EDA:} Kiến trúc Event-Driven Architecture đóng vai trò hấp thụ tải, trong khi Microservices mở rộng độc lập \cite{richards2020}.
    \item \textbf{Ưu tiên Energy Efficiency $\rightarrow$ EDA với giao thức bất đồng bộ:} Mô hình EDA loại bỏ hoàn toàn cơ chế lấy mẫu dữ liệu polling, tối đa hóa vòng đời nguồn cung cấp năng lượng \cite{lea2018}.
\end{itemize}

\vspace{1em}
\noindent\textbf{Ví dụ thực tiễn: Lựa chọn kiến trúc cho hệ thống bảo trì dự đoán điện gió ngoài khơi}

\textbf{1. Bối cảnh nghiệp vụ:} 
Một doanh nghiệp năng lượng cần triển khai hệ thống giám sát cho hàng trăm turbine ngoài khơi. Kênh giao tiếp duy nhất là mạng vệ tinh với băng thông hẹp, chi phí cao và độ trễ lớn \cite{ms_rollsroyce}.

\textbf{2. Đánh giá và ưu tiên Đặc tính kiến trúc:}
\begin{itemize}
    \item \textbf{Reliability:} Các quy trình an toàn phải vận hành tự chủ không phụ thuộc vào lệnh đám mây \cite{lea2018}.
    \item \textbf{Performance:} Việc phát hiện bất thường phải thực thi dưới 50ms, không thể chờ tín hiệu mạng vệ tinh \cite{lea2018}.
    \item \textbf{Scalability \& Bandwidth Efficiency:} Băng thông vệ tinh không cho phép truyền tải toàn bộ dữ liệu viễn trắc thô \cite{ms_rollsroyce}.
\end{itemize}

\textbf{3. Quyết định lựa chọn Mẫu kiến trúc:}
Mẫu \textbf{Edge-Cloud Hybrid} được chỉ định để giải quyết triệt để sự đánh đổi giữa 3 đặc tính trên \cite{lea2018}.

\textbf{4. Diễn giải giải pháp phân bổ:}
\begin{itemize}
    \item \textbf{Tại tầng Edge:} Thiết bị chạy Machine Learning cục bộ. Nếu phát hiện nguy hiểm, xuất lệnh ngắt với độ trễ < 10ms  \cite{lea2018}.
    \item \textbf{Tại tầng Cloud:} Tầng Edge chỉ gửi đi anomalies hoặc aggregated data, tiết kiệm 95\% băng thông. Cloud dùng để lưu trữ và huấn luyện lại Machine Learning  \cite{lea2018}.
\end{itemize}